% ═════════════════════════════════════════════════════════════════════════════
%  pages/abbreviations.tex
% ═════════════════════════════════════════════════════════════════════════════

% ─── Section 1 : Acronymes et termes techniques ───────────────────────────────
\chapter*{Acronymes et Termes Techniques}
\addcontentsline{toc}{chapter}{Acronymes et Termes Techniques}
\markboth{Acronymes et Termes Techniques}{Acronymes et Termes Techniques}

\begin{itemize}[leftmargin=2cm, labelwidth=1.8cm, align=left, itemsep=4pt]

  \item[\textbf{API}] Application Programming Interface (Interface de programmation applicative)

  \item[\textbf{CI/CD}] Continuous Integration / Continuous Deployment\\
        (Intégration Continue / Déploiement Continu)

  \item[\textbf{DAST}] Dynamic Application Security Testing\\
        (Test de sécurité dynamique des applications)

  \item[\textbf{DNS}] Domain Name System (Système de noms de domaine)

  \item[\textbf{gRPC}] Google Remote Procedure Call (Protocole d'appel de procédure à distance)

  \item[\textbf{IAST}] Interactive Application Security Testing\\
        (Test de sécurité interactif des applications)

  \item[\textbf{IDS}] Intrusion Detection System (Système de détection d'intrusion)

  \item[\textbf{JWT}] JSON Web Token (Jeton d'authentification web au format JSON)

  \item[\textbf{mTLS}] Mutual Transport Layer Security\\
        (Chiffrement TLS mutuel entre services)

  \item[\textbf{OWASP}] Open Worldwide Application Security Project

  \item[\textbf{RASP}] Runtime Application Self-Protection\\
        (Auto-protection applicative à l'exécution)

  \item[\textbf{RBAC}] Role-Based Access Control (Contrôle d'accès basé sur les rôles)

  \item[\textbf{SAST}] Static Application Security Testing\\
        (Test de sécurité statique des applications)

  \item[\textbf{STRIDE}] Spoofing, Tampering, Repudiation, Information Disclosure,\\
        Denial of Service, Elevation of Privilege\\
        (Méthodologie de modélisation des menaces)

  \item[\textbf{TLS}] Transport Layer Security (Protocole de sécurisation des communications)

  \item[\textbf{YAML}] YAML Ain't Markup Language (Format de sérialisation de données)

\end{itemize}

\newpage

% ─── Section 2 : Outils et technologies ──────────────────────────────────────
\chapter*{Outils et Technologies}
\addcontentsline{toc}{chapter}{Outils et Technologies}
\markboth{Outils et Technologies}{Outils et Technologies}

\begin{itemize}[leftmargin=3.2cm, labelwidth=3cm, align=left, itemsep=6pt]

  

  \item[\textbf{Envoy}] Proxy sidecar haute performance utilisé comme data plane dans
        Istio, responsable de l'interception et du traitement de tout le trafic interservices.

  \item[\textbf{Grafana}] Outil de visualisation et de création de tableaux de bord,
        utilisé pour l'observabilité des métriques de l'infrastructure.

  \item[\textbf{Istio}] Service Mesh open source déployé sur Kubernetes, fournissant
        le chiffrement mTLS, la gestion du trafic, l'authentification JWT et le contrôle
        d'accès RBAC de manière transparente.

  \item[\textbf{Kiali}] Console d'observabilité dédiée à Istio, offrant une visualisation
        graphique de la topologie des services et de leurs flux de communication.

  \item[\textbf{Kubernetes}] Plateforme d'orchestration de conteneurs permettant
        l'automatisation du déploiement, de la mise à l'échelle et de la gestion
        des applications conteneurisées.

  \item[\textbf{Minikube}] Outil permettant d'exécuter un cluster Kubernetes
        mono-nœud en local, utilisé dans le cadre du prototype de validation.


  \item[\textbf{Prometheus}] Système de collecte et de surveillance de métriques,
        intégré nativement à Istio pour l'observabilité du trafic interservices.

  

\end{itemize}